De bakker als vreetzak
droeg met veel ongemak
zijn nieuwbakken gebak

op zijn vlezig gestel.
Aan een leren bretel
zat een houten toestel,

die de taart die men vroeg,
direct voor zijn buik droeg.

Vraatzucht, zo'n vervelend trek
Waarom geven mensen zich eraan over?
Misschien omdat ze eenzaam zijn, of omdat ze gestresseerd
Of omdat ze gewoon niet weten hoe ze zich anders moeten bezighouden
Het kan ook zijn dat ze een ongezond voedingspatroon hebben
Of dat ze een tekort aan bepaalde voedingsstoffen hebben
Het is belangrijk om naar je lichaam te luisteren en gezonde keuzes te maken
Zodat je je goed voelt, zowel fysiek als mentaal
Laat vraatzucht achter je, kies voor een gezonde levensstijl!

zo groot als een
De bakker stak zijn piel in een stuk taart
en vroeg haar; “wat is je dit waard?”

Emmy had een enorme trek
maar nam niet deel aan het gesprek.

De vraag was immers retorisch,
Het was bedoeld als fallisch.

Zijn vrouw likte het zoete gebak.
en slikte daarna zijn zoute kwak.

De bakker legde uit; “als een man hijgt,
dan weet jij dat je zo wat krijgt.

Leer van teelvocht te leven,
want dat is wat we zullen geven!”

Er zijn diverse psychische redenen waarom mensen vraatzucht ontwikkelen. Hieronder staan een aantal van de meest voorkomende:

Stress: Mensen gebruiken vaak eten als een manier om te ontspannen of om stress te verminderen. Dit kan leiden tot een patroon van eetbuien of overmatig eten als reactie op stressvolle situaties.

Depressie: Depressie kan een vermindering van de eetlust veroorzaken, maar soms kan het ook leiden tot eetbuien of overmatig eten. Dit kan komen omdat mensen zich verdrietig of eenzaam voelen en eten als een manier om zich beter te voelen.

Angst: Angst kan leiden tot eetbuien of overmatig eten omdat mensen zich ongemakkelijk voelen en eten als een manier om zich beter te voelen.

Trauma: Mensen die een trauma hebben meegemaakt, zoals seksueel misbruik of emotioneel misbruik, kunnen een patroon van eetbuien of overmatig eten ontwikkelen als een manier om te cope met de emoties die het trauma oproept.

Perfectionisme: Mensen die perfectionistisch zijn, kunnen een patroon van eetbuien of overmatig eten ontwikkelen als een manier om de controle te behouden over hun lichaam en gewicht.

Er zijn diverse manieren om te werken aan vraatzucht. Hieronder staan enkele van de meest effectieve:

Cognitieve gedragstherapie (CGT): CGT richt zich op het veranderen van de gedachten en gedragspatronen die tot vraatzucht leiden. CGT kan helpen bij het on

Er zijn diverse manieren om te werken aan vraatzucht. Hieronder staan enkele van de meest effectieve:

Cognitieve gedragstherapie (CGT): CGT richt zich op het veranderen van de gedachten en gedragspatronen die tot vraatzucht leiden. CGT kan helpen bij het ontwikkelen van gezonde eetgewoonten en het verminderen van stress en emotionele eetbuien.

Gedragstherapie: Gedragstherapie richt zich op het veranderen van de gedragspatronen die tot vraatzucht leiden. Dit kan onder andere bestaan uit het bijhouden van een eetdagboek, het ontwikkelen van gezonde eetgewoonten en het aanleren van technieken om te dealen met eetbuien.

Interpersoonlijke psychotherapie: Interpersoonlijke psychotherapie richt zich op de relaties en interacties die een rol spelen in de vraatzucht. Dit kan helpen bij het verminderen van stress en emotionele eetbuien en het ontwikkelen van gezonde eetgewoonten.

Voedingsbegeleiding: Voedingsbegeleiding kan helpen bij het ontwikkelen van gezonde eetgewoonten en een gezonde relatie met eten. Een voedingsdeskundige kan helpen bij het opstellen van een gezonde maaltijdplan en het aanleren van technieken om gezonde keuzes te maken bij eten.

Medicatie: In sommige gevallen kan medicatie worden voorgeschreven om te helpen bij het verminderen van vraatzucht. Dit kan onder andere bestaan uit medicatie voor gewichtsverlies of medicatie die helpt bij het verminderen van emotionele eetbuien.

Het is belangrijk om te realiseren dat vraatzucht vaak een complexe aandoening is die ontstaat als gevolg van een combinatie van factoren. Het is daarom belangrijk om een behandelingsplan op maat te ontwikkelen dat rekening houdt met de unieke behoeften en omstandigheden van de persoon. Dit kan bestaan uit een combinatie van verschillende behandelingsmethoden, zoals therapeutische gesprekken, gedragstherapie, voedingsbegeleiding en medicatie.

Meditatie en mindfulness: Meditatie en mindfulness kunnen helpen bij het verminderen van stress en emotionele eetbuien en het ontwikkelen van een gezonde relatie met eten. Door middel van meditatie en mindfulness kunnen mensen leren om beter te luisteren naar hun lichaam en hun eetgedrag te observeren zonder oordeel. Dit kan helpen bij het identificeren van ongezonde eetpatronen en het ontwikkelen van gezonde alternatieven.

Als een taart wordt gegeten, gaat het eerst naar de mond, waar het wordt vermengd met speeksel. Hier begint de vertering al, enzymen in het speeksel beginnen het voedsel af te breken. Vanuit de mond gaat het naar de maag, waar het wordt vermengd met maagsap. Hier wordt het voedsel verder afgebroken door maagzuur en enzymen. Vanuit de maag gaat het naar de dunne darm, waar het voedsel verder wordt vertert en voedingsstoffen worden opgenomen in de bloedbaan.

Vanuit de dunne darm gaat het voedsel naar de dikke darm, waar overtollig water en elektrolyten worden opgenomen. Hier wordt het voedsel ook verder afgebroken door darmbacteriën. Vanuit de dikke darm gaat het naar de endeldarm, waar de ontlasting wordt gevormd. De ontlasting wordt dan naar de anus vervoerd en uitgescheiden uit het lichaam via de rectum.

Het is belangrijk om te weten dat de vertering van voedsel een complex proces is waarbij veel organen betrokken zijn. Dit proces kan variëren van persoon tot persoon, afhankelijk van hun gezondheid, dieet en andere factoren.

Als een taart wordt gegeten, worden er verschillende voedingsstoffen uit het taartdeeg en de vulling opgenomen in de bloedbaan. Hier zijn enkele van de belangrijkste voedingsstoffen die worden opgenomen:

Koolhydraten: De suikers en zetmeel in de taart worden afgebroken tot glucose, dat wordt opgenomen in de bloedbaan als energiebron voor de cellen in het lichaam.

Vetten: De vetten in de taart worden afgebroken in vetzuren, die worden opgenomen in de bloedbaan en gebruikt worden als energiebron of als bouwstof voor cellen.

Eiwitten: De eiwitten in de taart worden afgebroken in aminozuren, die worden opgenomen in de bloedbaan en gebruikt worden als bouwstof voor weefsels en hormonen.

Vitaminen en mineralen: De taart bevat ook kleine hoeveelheden vitaminen en mineralen, zoals vitamine A, vitamine C, vitamine D, ijzer, calcium, kalium, etc. Deze worden opgenomen in de bloedbaan en helpen bij de groei, de ontwikkeling en het onderhoud van het lichaam.

Het is belangrijk om op te merken dat taart vaak ook een aanzienlijke hoeveelheid calorieën en suiker bevat, dus het is aan te raden om met mate te eten. Overmatig eten van taart kan leiden tot overgewicht of andere gezondheidsproblemen.

Taart bevat vaak veel suiker en verzadigde vetten die niet goed zijn voor de gezondheid. Hier is een overzicht van enkele van de problemen die kunnen optreden wanneer je te veel taart eet:

Overgewicht en obesitas: Taart bevat vaak een aanzienlijke hoeveelheid calorieën en suiker, die kunnen bijdragen aan gewichtstoename als je teveel eet. Dit kan leiden tot overgewicht en obesitas, die verband houden met een verhoogd risico op hartziekten, diabetes en andere gezondheidsproblemen.

Tandproblemen: Taart bevat vaak veel suiker, die kan leiden tot tandbederf en andere tandproblemen.

Verhoogd cholesterol: Taart bevat vaak veel verzadigde vetten, die kunnen bijdragen aan een verhoogd cholesterolgehalte in het bloed. Dit kan leiden tot een verhoogd risico op hartziekten.

Leverproblemen: De verzadigde vetten in taart kunnen leiden tot een verhoogde vetaccumulatie in de lever, wat kan leiden tot leverproblemen zoals leververvetting.

Diabetes: Taart bevat vaak veel suiker, die kan leiden tot een verhoogd bloedsuikergehalte. Dit kan leiden tot een verhoogd risico op diabetes type 2.

Verzadigd vet wordt opgenomen in de bloedbaan via de dunne darm, na de vertering van de taart in de maag en de dunne darm. Eenmaal in de bloedbaan, wordt het verzadigde vet meegevoerd door lipoproteïnen zoals LDL (low-density lipoproteïnen) naar de lever. In de lever wordt het vet gebruikt voor energie of opgeslagen in vetcellen voor later gebruik.

Wanneer er teveel verzadigd vet in de bloedbaan aanwezig is, kan het leiden tot een verhoogd cholesterolgehalte. Dit komt doordat verzadigd vet de hoeveelheid LDL-cholesterol in het bloed verhoogt. LDL-cholesterol wordt vaak beschouwd als "slechte" cholesterol, omdat het zich kan ophopen in de bloedvaten en zo een verhoogd risico kan creëren voor hart- en vaatziekten.

Als er teveel LDL-cholesterol in het bloed aanwezig is, kan het zich ophopen in de wanden van de bloedvaten, vormen plaques die de bloedstroom belemmeren. Dit kan leiden tot een verhoogde bloeddruk en een verhoogd risico op hart- en vaatziekten zoals hartaanval en beroerte.

Ik ben een verzadigd vet molecuul,
in een taart geboren,
waar ik lag te rusten
in een zoete, romige oase.

Maar mijn reis was nog niet gedaan,
want na de vertering in de maag en dunne darm,
ging ik de bloedbaan in,
opgepakt door lipoproteïnen,
als een passagier in een trein.

Naar de lever gingen we,
waar ik werd opgeslagen,
of gebruikt voor energie,
maar als er teveel van mijn soort was,
kwam er een probleem in opspraak.

Want ik verhoogde het LDL-cholesterol,
en dat bracht risico's met zich mee.
Ik ophoopte in bloedvaten
en belemmerde de bloedstroom,
en dat was niet de bedoeling,
maar zo gaat het soms in de natuur.

Mijn reis eindigt hier niet,
want ik ben nog steeds in het lichaam,
maar als je gezond wilt blijven,
weet je nu, dat teveel van mij,
niet goed is voor het systeem.
Ik ben een suiker molecuul,
zoet van smaak en helder van kleur,
in een taart geboren,
waar ik lag te rusten
in een zoete, romige oase.

Maar mijn reis was nog niet gedaan,
want na de vertering in de mond,
ging ik de bloedbaan in,
waar ik werd opgenomen als glucose,
de brandstof voor de cellen in het lichaam.

Maar als ik teveel werd geconsumeerd,
kwam er een probleem in opspraak.
Ik verhoogde het bloedsuikergehalte,
en dat bracht risico's met zich mee,
zoals diabetes en overgewicht.

Mijn reis eindigt hier niet,
want ik ben nog steeds in het lichaam,
maar als je gezond wilt blijven,
weet je nu, dat teveel van mij,
niet goed is voor het systeem.

Ik ben een suiker molecuul,
zoet van smaak en helder van kleur,
maar in te grote hoeveelheden,
kan ik problemen veroorzaken,
zoals diabetes, de ziekte van de suiker.

Diabetes is een aandoening waarbij
de cellen geen glucose meer kunnen opnemen
om te gebruiken als energie,
omdat het lichaam geen insuline meer aanmaakt of deze niet goed gebruikt.

Ik blijf circuleren in het bloed,
waardoor de bloedsuikerspiegel stijgt,
en de cellen geen energie hebben,
waardoor het lichaam uitgeput raakt.

Diabetes is een ernstige ziekte,
die kan leiden tot ernstige complicaties,
als het niet goed wordt behandeld,
maar met een gezonde levensstijl en medicatie,
kan het worden gecontroleerd.

Dus als je me wilt genieten,
doet dat met mate,
want teveel van mij,
kan leiden tot diabetes,
de ziekte van de suiker
